\documentclass[12pt, a4paper]{article}
\usepackage{../../../myconfig} % Gọi file cấu hình từ ngoài gốc vào

% ==========================================
% BẮT ĐẦU NỘI DUNG TÀI LIỆU
% ==========================================
\begin{document}

% Truyền 3 tham số: {Nhãn cam} {Chữ to chính giữa} {Chữ ở góc phải Header}
\titlebox{Chủ đề: Tích phân}{Nguyên hàm}{Nguyên hàm - Tích phân: Nguyên hàm}
\drawdashlines{20}
\newpage
\drawdashlines{25}
\newpage
\drawdashlines{25}
\newpage
\drawdashlines{25}
\newpage

% --- CÂU 1 ---
\questionLinesManual{7}{1}{Hàm số \( F(x) \) là một nguyên hàm của hàm số \( f(x) \) trên khoảng \( K \) nếu}{
    \begin{tasks}(2)
        \task \( F'(x) = -f(x), \forall x \in K \)
        \task \( f'(x) = F(x), \forall x \in K \)
        \task \( F'(x) = f(x), \forall x \in K \)
        \task \( f'(x) = -F(x), \forall x \in K \)
    \end{tasks}
}

% --- CÂU 2 ---
\questionLinesManual{7}{2}{Họ nguyên hàm của hàm số \( f(x) = x^2 \) là}{
    \begin{tasks}(4)
        \task \( 3x^2 + C \)
        \task \( 2x^3 + C \)
        \task \( \dfrac{1}{3}x^3 + C \)
        \task \( \dfrac{1}{2}x^3 + C \)
    \end{tasks}
}

% --- CÂU 3 ---
\questionLinesManual{7}{3}{Một nguyên hàm \( F(x) \) của hàm số \( f(x) = 3x^2 - 2x + 2025 \) là}{
    \begin{tasks}(4)
        \task \( F(x) = x^3 - x^2 + 2025x \)
        \task \( F(x) = x^3 - x^2 + 2025 \)
        \task \( F(x) = 6x - 2 \)
        \task \( F(x) = x^3 - x^2 \)
    \end{tasks}
}

% --- CÂU 4 ---
\questionLinesManual{10}{4}{Họ các nguyên hàm của hàm số \( f(x) = \dfrac{1}{x^3} \) là}{
    \begin{tasks}(4)
        \task \( -\dfrac{3}{x^4} + C \)
        \task \( -\dfrac{1}{x^2} + C \)
        \task \( -\dfrac{1}{2x^2} + C \)
        \task \( -\dfrac{1}{4x^4} + C \)
    \end{tasks}
}

% --- CÂU 6 ---
\questionLinesManual{10}{5}{Họ nguyên hàm của hàm số \( f(x) = \dfrac{3}{x^2} \) là}

% --- CÂU 5 ---
\questionLinesManual{8}{6}{Hàm số nào sau đây \textbf{không} là một nguyên hàm của \( f(x) = \sqrt[3]{x} \) trên \( (0; +\infty) \)?}{
    \begin{tasks}(4)
        \task \( F_1(x) = \dfrac{3\sqrt[3]{x^4}}{4} + 1 \)
        \task \( F_3(x) = \dfrac{3x\sqrt[3]{x}}{4} + 3 \)
        \task \( F_4(x) = \dfrac{3}{4}x^{\frac{4}{3}} + 4 \)
        \task \( F_2(x) = \dfrac{3\sqrt[4]{x^3}}{4} + 2 \)
    \end{tasks}
}

% --- CÂU 7 ---
\questionLinesManual{8}{7}{Họ nguyên hàm của hàm số \( f(x) = \dfrac{1}{\sqrt[3]{x}} \) với \( x > 0 \) là}

% --- CÂU 8 ---
\questionLinesManual{8}{8}{Tìm nguyên hàm của hàm số \( f(x) = \dfrac{1}{x - 2} \).}

% --- CÂU 9 ---
\questionLinesManual{8}{9}{Tìm họ nguyên hàm của hàm số \( f(x) = \dfrac{1}{2x+3} \).}{
    \begin{tasks}(4)
        \task \( \ln|2x+3| + C \)
        \task \( \dfrac{1}{2} \ln|2x+3| + C \)
        \task \( \dfrac{1}{\ln 2} \ln|2x+3| + C \)
        \task \( \dfrac{1}{2} \ln(2x+3) + C \)
    \end{tasks}
}

% --- CÂU 10 ---
\questionLinesManual{6}{10}{Họ nguyên hàm của hàm số \( f(x) = \sin x + \cos x \) là}{
    \begin{tasks}(4)
        \task \( \cos x + \sin x + C \)
        \task \( \cos x - \sin x + C \)
        \task \( -\cos x - \sin x + C \)
        \task \( -\cos x + \sin x + C \)
    \end{tasks}
}

% --- CÂU 11 ---
\questionLinesManual{6}{11}{\( \displaystyle \int (\sin x + 4x^3) \, dx \) bằng}{
    \begin{tasks}(4)
        \task \( -\cos x + 4x^4 + C \)
        \task \( \cos x + x^4 + C \)
        \task \( \cos x + 12x^2 + C \)
        \task \( -\cos x + x^4 + C \)
    \end{tasks}
}

% --- CÂU 12 ---
\questionLinesManual{10}{12}{Tìm nguyên hàm của hàm số \( f(x) = \cos 3x \).}{
    \begin{tasks}(2)
        \task \( \displaystyle \int \cos 3x \, dx = 3\sin 3x + C \)
        \task \( \displaystyle \int \cos 3x \, dx = \dfrac{\sin 3x}{3} + C \)
        \task \( \displaystyle \int \cos 3x \, dx = \sin 3x + C \)
        \task \( \displaystyle \int \cos 3x \, dx = -\dfrac{\sin 3x}{3} + C \)
    \end{tasks}
}

% --- CÂU 13 ---
\questionLinesManual{10}{13}{Tìm nguyên hàm của hàm số \( f(x) = 7^x \).}{
    \begin{tasks}(2)
        \task \( \displaystyle \int 7^x \, dx = \dfrac{7^x}{\ln 7} + C \)
        \task \( \displaystyle \int 7^x \, dx = 7^{x+1} + C \)
        \task \( \displaystyle \int 7^x \, dx = \dfrac{7^{x+1}}{x+1} + C \)
        \task \( \displaystyle \int 7^x \, dx = 7^x \ln 7 + C \)
    \end{tasks}
}

% --- CÂU 14 ---
\questionLinesManual{8}{14}{Họ nguyên hàm của hàm số \( f(x) = e^x + x \) là}{
    \begin{tasks}(4)
        \task \( e^x + 1 + C \)
        \task \( e^x + x^2 + C \)
        \task \( e^x + \dfrac{1}{2}x^2 + C \)
        \task \( \dfrac{1}{x+1}e^x + \dfrac{1}{2}x^2 + C \)
    \end{tasks}
}

% --- CÂU 15 ---
\questionLinesManual{8}{15}{Họ nguyên hàm của hàm số \( f(x) = e^{3x} \) là hàm số nào sau đây?}{
    \begin{tasks}(4)
        \task \( 3e^x + C \)
        \task \( \dfrac{1}{3}e^{3x} + C \)
        \task \( \dfrac{1}{3}e^x + C \)
        \task \( 3e^{3x} + C \)
    \end{tasks}
}

% --- CÂU 16 ---
\questionLinesManual{8}{16}{Tìm nguyên hàm của hàm số \( f(x) = \dfrac{x^4 + 2}{x^2} \).}

% --- CÂU 17 ---
\questionLinesManual{12}{17}{Cho hàm số \( f(x) = \dfrac{2x^4 + 3}{x^2} \). Tìm nguyên hàm của hàm số đã cho.}

% --- CÂU 18 ---
\questionLinesManual{12}{18}{Cho \( y = F(x) \) là một nguyên hàm của hàm số \( f(x) = -3x^2 + 4x + 2 \) và \( F(1) = 2 \). Tính \( F(-1) \).}{
    \begin{tasks}(4)
        \task \( F(-1) = 0 \)
        \task \( F(-1) = 4 \)
        \task \( F(-1) = 5 \)
        \task \( F(-1) = 1 \)
    \end{tasks}
}

% --- CÂU 19 ---
\questionLinesManual{8}{19}{Cho \( F(x) \) là một nguyên hàm của hàm số \( f(x) = e^x + 2x \) thỏa mãn \( F(0) = 1 \). Khẳng định nào sau đây là đúng?}{
    \begin{tasks}(4)
        \task \( F(x) = e^x + x^2 + 1 \)
        \task \( F(x) = e^x + x + 1 \)
        \task \( F(x) = e^x + x^2 \)
        \task \( F(x) = e^x + 2x^2 \)
    \end{tasks}
}

% --- CÂU 20 ---
\questionLinesManual{12}{20}{Cho hàm số \( f(x) \) xác định trên \( \mathbb{R} \setminus \left\{ \dfrac{1}{2} \right\} \) thỏa mãn \(f'(x) = \frac{2}{2x-1}, \, f(0) = 1, \, f(1) = 2.\) Giá trị của biểu thức \( f(-1) + f(3) \) bằng}{
    \begin{tasks}(4)
        \task \( 2 + \ln 15 \)
        \task \( 3 + \ln 15 \)
        \task \( \ln 15 \)
        \task \( 4 + \ln 15 \)
    \end{tasks}
}
\newpage
\drawdashlines{25}

\end{document}
